% !TEX root = ../main.tex
\chapter{Finite-sample study}
\label{ch:finite-sample}

This chapter reports a finite-sample study for the estimator developed in
\Cref{ch:pooled-extreme}. The purpose is deliberately narrow. We study the
pooled extreme-expectile estimator
\[
  \widehat\xi_{\tau}^{\mathrm{pool},\star}
    (\bm\nu_n^k,\bm\omega)
  =
  \psifn\!\big(\hat\gamma_n(\bm\nu_n^k)\big)
  \hatWei_n(1-p\mid\bm\omega),
  \qquad
  \bm\nu_n^k=(k_1/k,\ldots,k_m/k)^\top,
\]
where \(p=1-\tau\) and \(k=\sum_{j=1}^{m}k_j\). Thus the bridge-Hill
weights are fixed at the deterministic threshold-allocation convention,
whereas the pooled-Weissman weights \(\bm\omega\) vary across the
Daouia--Padoan--Stupfler choices studied below. The simulations do not claim
optimality of \(\bm\nu_n^k\), do not use random bridge-Hill weights, and do
not revive the intermediate expectile-pooling route discarded earlier.

The main empirical message is also narrow. In the representative imbalanced
design, DPS variance weights make the distributed estimator behave almost like
the centralised benchmark for point estimation, while equal weights are mildly
less efficient and AMSE weights add no systematic gain. The fixed bridge-Hill
convention \(\bm\nu_n^k\) is not visibly driving the conclusions. The main
limitation is interval calibration: the first-order log intervals often
under-cover, and the decomposition diagnostics suggest that finite-sample
bridge effects and lower-order Hill effects can still be visible even though
they disappear in the theorem.

\section{Estimators and implementation}
\label{sec:sim:estimators}

For each machine \(j\), the local Hill estimator and Weissman quantile
estimator are computed using the \texttt{HTailIndex} and \texttt{extQuantile}
routines from the \texttt{ExtremeRisks} package \parencite{ExtremeRisksR}.
The simulation code only adds the thesis-specific pooling layer: geometric
pooling of the local Weissman estimates, the map \(\psifn\), the deterministic
bridge-Hill convention \(\bm\nu_n^k\), and the summary calculations.

The following estimators are recorded in every scenario:
\begin{itemize}
\item \emph{equal}: \(\omega_j=1/m\);
\item \emph{DPS variance}: \(\omega_j=k_j/k\), the finite-sample independent
      distributed version of the DPS variance weights;
\item \emph{DPS AMSE oracle}: the DPS AMSE formula evaluated with the
      population second-order parameters of the data-generating model;
\item \emph{DPS AMSE plug-in}: the same formula with second-order inputs
      estimated using \texttt{evt0} \parencite{evt0R};
\item \emph{centralised}: the corresponding single-sample estimator computed
      from the combined sample, used only as a benchmark.
\end{itemize}
All distributed estimators use the same bridge-Hill vector
\(\bm\nu_n^k\). In particular, the plug-in AMSE estimator randomises
\(\bm\omega\) only.

For each estimator, inference is assessed with the log-scale interval from
\Cref{cor:pool-ext:plugin-inference}. In a Monte Carlo replication, if
\(\widehat B_\omega\) and \(\widehat V_\omega\) denote the selected DPS
centring and variance objects, the interval is
\[
  \widehat\xi
  \left[
    \exp\!\left\{-\frac{\ell}{\sqrt{k}}
      \big(\widehat B_\omega+z_{0.975}\sqrt{\widehat V_\omega}\big)\right\},
    \exp\!\left\{-\frac{\ell}{\sqrt{k}}
      \big(\widehat B_\omega-z_{0.975}\sqrt{\widehat V_\omega}\big)\right\}
  \right],
  \qquad
  \ell=\log\!\left(\frac{k}{np}\right).
\]
The interval is included as a diagnostic of the first-order plug-in route,
not as a claim of finite-sample calibration.

\section{Design}
\label{sec:sim:design}

The study stays within the iid common-marginal setting of the theorem. The
data-generating models are Pareto with \(\gamma=0.25\), Pareto with
\(\gamma=0.60\), Burr with \(\gamma=0.25\) and \(\rho=-1\), Burr with
\(\gamma=0.60\) and \(\rho=-1\), and the absolute Student distribution with
three degrees of freedom. The latter has \(\gamma=1/3\). The true expectile
\(\xi_\tau\) is computed numerically from the population expectile equation
\[
  \tau\,\E[(X-\xi_\tau)_+]
  =
  (1-\tau)\,\E[(\xi_\tau-X)_+],
\]
rather than by the asymptotic expectile--quantile bridge.

The final simulation uses \(5000\) Monte Carlo replications per scenario. The
main design has \(n=10000\), \(k_j/n_j=0.05\), \(np=5\),
\(m\in\{5,10\}\), and either balanced allocation or a strong imbalanced
allocation with machine proportions increasing geometrically from \(1\) to
\(10\). Two sensitivity blocks then fix \(m=10\) and the strong allocation:
one varies \(k_j/n_j\in\{0.03,0.05,0.10\}\), and the other varies
\(np\in\{1,5,10\}\).

Performance is measured by the Monte Carlo mean of the log relative error,
the root mean squared log relative error, empirical coverage of the nominal
\(95\%\) log interval, average log interval width, invalid-estimate rate, and
weight-stability diagnostics. The code also records the exact finite-sample
decomposition terms
\[
  A_n(\bm\omega)
  =
  \log\{\hatWei_n(1-p\mid\bm\omega)/Q(1-p)\},
  \qquad
  B_n(\bm\nu)
  =
  \log\psifn\{\hat\gamma_n(\bm\nu)\}-\log\psifn(\gamma),
\]
and
\[
  C_n
  =
  \log\{\xi_\tau/(\psifn(\gamma)Q(1-p))\}.
\]
These diagnostics are included to connect the Monte Carlo evidence to
\Cref{lem:pool-ext:decomposition}; they are not used to estimate new
asymptotic constants.

\section{Point estimation and pooled-Weissman weights}
\label{sec:sim:point}

\Cref{tab:sim:point-main} reports the representative main design with
\(m=10\), strong sample-size imbalance, \(k_j/n_j=0.05\), and \(np=5\). The
table compares the simple equal weights, the DPS variance weights, the oracle
AMSE weights, the plug-in AMSE weights, and the centralised benchmark. The
main result is that distributed pooling loses very little relative to the
centralised benchmark when the pooled-Weissman weights are chosen by the DPS
variance rule. The log RMSEs are identical to three decimals in the Burr heavy
case, essentially identical in the Pareto heavy and Student cases, and only
slightly larger in the light-tailed Pareto and Burr cases. This is the main
positive finite-sample evidence for the promoted estimator: the communication
restriction does not create a large point-estimation penalty in this grid.

The comparison across distributed weights gives a simple practical lesson.
Equal weights are consistently worse under strong imbalance, especially in the
heavier-tailed models, because they ignore the different local sample sizes.
The oracle and plug-in AMSE weights are almost indistinguishable from the
variance weights in the representative design. Thus the simulation supports
DPS variance weights as a stable practical baseline for \(\bm\omega\), but it
does not show a finite-sample benefit from estimating second-order AMSE
weights here. This is a statement about the pooled-Weissman weights only; it
does not suggest opening a new optimisation layer for \(\bm\nu\).

\begin{table}[htbp]
\centering
\small
\caption{Point-estimation performance in the representative main design
\((m=10\), strong allocation, \(k_j/n_j=0.05\), \(np=5)\).}
\label{tab:sim:point-main}
\begingroup
\setlength{\tabcolsep}{4pt}
\begin{tabular}{lrllrrr}
\toprule
DGP & $m$ & Regime & Estimator & Log bias & Log RMSE & Valid rate \\
\midrule
Burr heavy & 10 & strong & Centralised & 0.077 (0.003) & 0.226 (0.002) & 1.000 (0.000) \\
Burr heavy & 10 & strong & DPS AMSE oracle & 0.075 (0.003) & 0.223 (0.002) & 1.000 (0.000) \\
Burr heavy & 10 & strong & DPS AMSE plug-in & 0.075 (0.003) & 0.223 (0.002) & 1.000 (0.000) \\
Burr heavy & 10 & strong & DPS variance & 0.075 (0.003) & 0.223 (0.002) & 1.000 (0.000) \\
Burr heavy & 10 & strong & Equal & 0.070 (0.003) & 0.245 (0.003) & 1.000 (0.000) \\
Burr light & 10 & strong & Centralised & -0.039 (0.001) & 0.069 (0.001) & 1.000 (0.000) \\
Burr light & 10 & strong & DPS AMSE oracle & -0.041 (0.001) & 0.070 (0.001) & 1.000 (0.000) \\
Burr light & 10 & strong & DPS AMSE plug-in & -0.041 (0.001) & 0.070 (0.001) & 1.000 (0.000) \\
Burr light & 10 & strong & DPS variance & -0.041 (0.001) & 0.070 (0.001) & 1.000 (0.000) \\
Burr light & 10 & strong & Equal & -0.043 (0.001) & 0.084 (0.001) & 1.000 (0.000) \\
Pareto heavy & 10 & strong & Centralised & -0.007 (0.003) & 0.204 (0.002) & 1.000 (0.000) \\
Pareto heavy & 10 & strong & DPS AMSE oracle & -0.012 (0.003) & 0.204 (0.002) & 1.000 (0.000) \\
Pareto heavy & 10 & strong & DPS AMSE plug-in & -- & -- & -- \\
Pareto heavy & 10 & strong & DPS variance & -0.012 (0.003) & 0.204 (0.002) & 1.000 (0.000) \\
Pareto heavy & 10 & strong & Equal & -0.012 (0.003) & 0.231 (0.002) & 1.000 (0.000) \\
Pareto light & 10 & strong & Centralised & -0.069 (0.001) & 0.089 (0.001) & 1.000 (0.000) \\
Pareto light & 10 & strong & DPS AMSE oracle & -0.071 (0.001) & 0.090 (0.001) & 1.000 (0.000) \\
Pareto light & 10 & strong & DPS AMSE plug-in & -- & -- & -- \\
Pareto light & 10 & strong & DPS variance & -0.071 (0.001) & 0.090 (0.001) & 1.000 (0.000) \\
Pareto light & 10 & strong & Equal & -0.071 (0.001) & 0.100 (0.001) & 1.000 (0.000) \\
Abs. Student-3 & 10 & strong & Centralised & 0.110 (0.001) & 0.143 (0.001) & 1.000 (0.000) \\
Abs. Student-3 & 10 & strong & DPS AMSE oracle & 0.109 (0.001) & 0.142 (0.001) & 1.000 (0.000) \\
Abs. Student-3 & 10 & strong & DPS AMSE plug-in & 0.109 (0.001) & 0.142 (0.001) & 1.000 (0.000) \\
Abs. Student-3 & 10 & strong & DPS variance & 0.109 (0.001) & 0.142 (0.001) & 1.000 (0.000) \\
Abs. Student-3 & 10 & strong & Equal & 0.108 (0.002) & 0.154 (0.001) & 1.000 (0.000) \\
\bottomrule
\end{tabular}

\endgroup
\end{table}

\begin{figure}[htbp]
\centering
\includegraphics[width=\textwidth]{generated/simulation/figures/rmse_pooled_vs_centralised.pdf}
\caption{Log RMSE of the DPS variance pooled estimator against the
centralised benchmark in the main grid with \(k_j/n_j=0.05\) and \(np=5\).
The dashed line is equality. Points near the line indicate that the pooled
route tracks the centralised benchmark closely.}
\label{fig:sim:rmse-estimator}
\end{figure}

\section{Deterministic bridge-Hill sensitivity}
\label{sec:sim:nu-sensitivity}

The main specification fixes the bridge-Hill weights at
\(\bm\nu_n^k=(k_1/k,\ldots,k_m/k)^\top\). To check whether this convention has
visible finite-sample consequences, \Cref{tab:sim:nu} repeats the representative
main design for the DPS variance pooled-Weissman weights while replacing
\(\bm\nu_n^k\) by two other deterministic conventions: equal bridge-Hill
weights and sample-allocation weights \((n_1/n,\ldots,n_m/n)^\top\). This is
only a sensitivity check. It does not search for an optimal \(\bm\nu\), and it
does not use the random estimated \(\bm\omega\) as a bridge-Hill vector.

The rows are close. In this design the threshold and sample-allocation
conventions coincide because the same \(k_j/n_j\) ratio is used on every
machine, and their reported performance is identical. Replacing them by equal
bridge-Hill weights changes little in the light-tailed Burr and Pareto cases
and in the Student case. It has a modest cost in the heavier-tailed designs:
for example, log RMSE increases from \(0.223\) to \(0.236\) in Burr heavy and
from \(0.204\) to \(0.216\) in Pareto heavy, with slightly lower coverage.
The practical conclusion is that \(\bm\nu_n^k\) is not driving the main
finite-sample conclusions, although the bridge-Hill component can still matter
at a lower-order level in heavier-tailed settings.

\begin{table}[htbp]
\centering
\small
\caption{Sensitivity to deterministic bridge-Hill weights in the representative
main design. The pooled-Weissman weights are fixed at the DPS variance
choice.}
\label{tab:sim:nu}
\begingroup
\setlength{\tabcolsep}{4pt}
\begin{tabular}{llrrr}
\toprule
DGP & Bridge weights & Log bias & Log RMSE & Coverage \\
\midrule
Burr heavy & Equal & 0.077 (0.003) & 0.236 (0.003) & 0.720 (0.006) \\
Burr heavy & Sample & 0.075 (0.003) & 0.223 (0.002) & 0.751 (0.006) \\
Burr heavy & $\nu^k$ & 0.075 (0.003) & 0.223 (0.002) & 0.751 (0.006) \\
Burr light & Equal & -0.041 (0.001) & 0.070 (0.001) & 0.593 (0.007) \\
Burr light & Sample & -0.041 (0.001) & 0.070 (0.001) & 0.591 (0.007) \\
Burr light & $\nu^k$ & -0.041 (0.001) & 0.070 (0.001) & 0.591 (0.007) \\
Pareto heavy & Equal & -0.007 (0.003) & 0.216 (0.002) & -- \\
Pareto heavy & Sample & -0.012 (0.003) & 0.204 (0.002) & -- \\
Pareto heavy & $\nu^k$ & -0.012 (0.003) & 0.204 (0.002) & -- \\
Pareto light & Equal & -0.071 (0.001) & 0.090 (0.001) & -- \\
Pareto light & Sample & -0.071 (0.001) & 0.090 (0.001) & -- \\
Pareto light & $\nu^k$ & -0.071 (0.001) & 0.090 (0.001) & -- \\
Abs. Student-3 & Equal & 0.110 (0.001) & 0.144 (0.001) & 0.910 (0.004) \\
Abs. Student-3 & Sample & 0.109 (0.001) & 0.142 (0.001) & 0.920 (0.004) \\
Abs. Student-3 & $\nu^k$ & 0.109 (0.001) & 0.142 (0.001) & 0.920 (0.004) \\
\bottomrule
\end{tabular}

\endgroup
\end{table}

\section{Plug-in interval performance}
\label{sec:sim:intervals}

\Cref{tab:sim:interval-main} shows that the first-order log intervals are
not uniformly well calibrated at this sample size. Coverage is close to
nominal only in the absolute Student case, where the DPS variance and plug-in
AMSE intervals cover about \(92\%\) of replications. In the Pareto and Burr
cases, coverage is materially below \(95\%\), ranging from about \(59\%\) in
the light Burr case to about \(75\%\) in the heavy Burr case. The centralised
benchmark also under-covers in every model, so the problem is not simply that
the sample has been distributed across machines. At the same time, the size of
the discrepancy is model-dependent: distributed pooling is much worse than
the centralised benchmark in the light Burr case, while it is much better in
the Student case. The most defensible reading is therefore that interval
calibration is a finite-sample weakness of the first-order plug-in route, with
both bridge and centring effects and the distributed construction contributing
differently across models.

The interval results therefore support a cautious interpretation of
\Cref{cor:pool-ext:plugin-inference}. The corollary gives the asymptotic
studentisation when the DPS source route supplies consistent plug-in objects;
the finite-sample study shows that the resulting interval can still be
noticeably under-calibrated in moderate samples and biased tail models.

\Cref{fig:sim:bridge-diagnostic} gives a targeted diagnostic for this
undercoverage. It evaluates the same oracle first-order interval against both
the exact population expectile and the first-order bridge target
\(\psifn(\gamma)Q(1-p)\). In the light Pareto and Burr designs, much of the
coverage loss against the exact expectile disappears when the bridge target is
used, and exact-target coverage improves in the larger-\(n\), smaller
\(k/n\) diagnostic design. This supports the interpretation that the main
undercoverage is a plausible finite-sample bridge and first-order calibration
effect, not a sign error or a failure of the pooling construction. The
diagnostic is not uniform across all heavy-tailed designs, so the conclusion
remains cautionary rather than corrective.

\begin{table}[htbp]
\centering
\small
\caption{Empirical coverage and log-width of nominal \(95\%\) intervals in
the representative main design.}
\label{tab:sim:interval-main}
\begingroup
\setlength{\tabcolsep}{4pt}
\begin{tabular}{lrllrrrrr}
\toprule
DGP & $m$ & Regime & Estimator & Coverage & Log width & CI rate & $\eta_n$ & Bridge gap \\
\midrule
Burr heavy & 10 & strong & Centralised & 0.400 & 0.744 & 1.000 & 0.083 & 0.021 \\
Burr heavy & 10 & strong & DPS AMSE plug-in & 0.300 & 0.743 & 1.000 & 0.083 & 0.021 \\
Burr heavy & 10 & strong & DPS variance & 0.400 & 0.743 & 1.000 & 0.083 & 0.021 \\
\bottomrule
\end{tabular}

\endgroup
\end{table}

\begin{figure}[htbp]
\centering
\includegraphics[width=\textwidth]{generated/simulation/figures/coverage_heatmap.pdf}
\caption{Empirical coverage in the representative main design
\((m=10\), strong allocation, \(k_j/n_j=0.05\), \(np=5)\). Cell labels give
coverage; colour records the deviation from the nominal \(95\%\) level.}
\label{fig:sim:coverage-estimator}
\end{figure}

\begin{figure}[htbp]
\centering
\includegraphics[width=\textwidth]{generated/simulation/figures/bridge_target_diagnostic.pdf}
\caption{Interval diagnostic comparing empirical coverage against the exact
population expectile and against the first-order bridge target
\(\psifn(\gamma)Q(1-p)\). The diagnostic uses oracle DPS centring and
variance objects and is separate from the final main grid.}
\label{fig:sim:bridge-diagnostic}
\end{figure}

\section{Decomposition diagnostics}
\label{sec:sim:decomposition}

\Cref{tab:sim:decomposition} records the scaled decomposition terms in the
representative main design for the DPS variance estimator with
\(\bm\nu=\bm\nu_n^k\). The columns report root mean squared values of
\((\sqrt{k}/\ell)A_n\), \((\sqrt{k}/\ell)B_n\), and
\((\sqrt{k}/\ell)C_n\), together with the scaled log error. The final column is
the numerical reconstruction check for
\[
  \log(\widehat\xi_\tau/\xi_\tau)=A_n(\bm\omega)+B_n(\bm\nu)-C_n.
\]

This table is meant to make the simulation chapter answer the same question as
the proof: which terms are visible at the pooled-Weissman scale? In the
asymptotic theorem \(A_n(\bm\omega)\) carries the first-order law, while
\(B_n(\bm\nu)\) and \(C_n\) vanish after multiplication by
\(\sqrt{k}/\ell\). At the simulated sample size, that separation is only
partly visible. The pooled-Weissman component remains the largest stochastic
term in most rows, but the bridge-Hill term is not negligible in the
heavier-tailed cases: its RMS scaled value is \(0.474\) in Burr heavy and
\(0.382\) in Pareto heavy, compared with \(0.638\) and \(0.610\) for
\(A_n(\bm\omega)\). In contrast, \(B_n(\bm\nu)\) is almost invisible in the
light-tailed Pareto and Burr designs.

The population bridge term \(C_n\) shows the opposite pattern. It is small in
the heavier-tailed Pareto and Burr cases, but it is comparable to, or larger
than, the pooled-Weissman component in the light-tailed cases. This helps
explain why point-estimation performance and interval calibration need not
move together: some finite-sample error is coming from the population
expectile--quantile bridge rather than from the distributed quantile pooling
itself. The zero reconstruction errors in the final column confirm that these
diagnostics are reading the exact decomposition rather than an approximate
post-hoc split.

\begin{table}[htbp]
\centering
\small
\caption{Finite-sample decomposition ledger in the representative main design.
The estimator uses DPS variance pooled-Weissman weights and
\(\bm\nu=\bm\nu_n^k\).}
\label{tab:sim:decomposition}
\begingroup
\setlength{\tabcolsep}{4pt}
\begin{tabular}{lrrrrr}
\toprule
DGP & RMS scaled log error & RMS scaled $A_n$ & RMS scaled $B_n$ & $|$scaled $C_n|$ & Check \\
\midrule
Burr heavy & 1.080 & 0.638 & 0.474 & 0.045 & 0.00 \\
Burr light & 0.340 & 0.269 & 0.019 & 0.279 & 0.00 \\
Pareto heavy & 0.987 & 0.610 & 0.382 & 0.057 & 0.00 \\
Pareto light & 0.437 & 0.255 & 0.013 & 0.338 & 0.00 \\
Abs. Student-3 & 0.690 & 0.589 & 0.222 & 0.143 & 0.00 \\
\bottomrule
\end{tabular}

\endgroup
\end{table}

\section{Sensitivity to thresholds and targets}
\label{sec:sim:sensitivity}

The threshold sensitivity in \Cref{tab:sim:threshold} shows the usual
bias--variance trade-off, but the important point is that RMSE and coverage
often point in different directions. For Burr light, increasing \(k_j/n_j\)
from \(0.03\) to \(0.10\) lowers the DPS variance log RMSE from \(0.084\) to
\(0.047\), while coverage falls from about \(70\%\) to \(31\%\). For Pareto
heavy, the larger threshold improves RMSE, consistent with the exact Pareto
model having no second-order bias of the Burr type. For the absolute Student
case, the largest threshold worsens RMSE even though coverage remains around
\(92\%\). Thus threshold choice cannot be justified by point error alone.

\Cref{tab:sim:target} varies the target through \(np\). The least extreme
target, \(np=10\), often has smaller log RMSE than \(np=1\), as expected from
the shorter extrapolation distance. Coverage, however, usually falls as
\(np\) increases in this grid. The pattern is especially sharp in the light
Burr and light Pareto cases, where DPS variance coverage drops from about
\(78\%\) and \(84\%\), respectively, at \(np=1\), to about \(43\%\) and
\(49\%\) at \(np=10\). This again points to the interaction between
second-order bias, threshold choice, and first-order interval centring, rather
than to a simple monotone improvement when the target is made less extreme.

\begin{table}[htbp]
\centering
\small
\caption{Threshold sensitivity for \(m=10\), strong allocation and \(np=5\).}
\label{tab:sim:threshold}
\begingroup
\setlength{\tabcolsep}{4pt}
\begin{tabular}{lllrrr}
\toprule
DGP & Threshold rule & Estimator & $k/n$ & Log RMSE & Coverage \\
\midrule
Burr heavy & fraction=0.03 & DPS AMSE plug-in & 0.029 & 0.255 & 0.731 \\
Burr heavy & fraction=0.03 & DPS variance & 0.029 & 0.255 & 0.731 \\
Burr heavy & fraction=0.05 & DPS AMSE plug-in & 0.050 & 0.228 & 0.732 \\
Burr heavy & fraction=0.05 & DPS variance & 0.050 & 0.228 & 0.733 \\
Burr heavy & fraction=0.1 & DPS AMSE plug-in & 0.100 & 0.260 & 0.744 \\
Burr heavy & fraction=0.1 & DPS variance & 0.100 & 0.260 & 0.743 \\
Burr light & fraction=0.03 & DPS AMSE plug-in & 0.029 & 0.084 & 0.695 \\
Burr light & fraction=0.03 & DPS variance & 0.029 & 0.084 & 0.696 \\
Burr light & fraction=0.05 & DPS AMSE plug-in & 0.050 & 0.069 & 0.580 \\
Burr light & fraction=0.05 & DPS variance & 0.050 & 0.069 & 0.579 \\
Burr light & fraction=0.1 & DPS AMSE plug-in & 0.100 & 0.047 & 0.314 \\
Burr light & fraction=0.1 & DPS variance & 0.100 & 0.047 & 0.316 \\
Pareto heavy & fraction=0.03 & DPS AMSE plug-in & 0.029 & -- & -- \\
Pareto heavy & fraction=0.03 & DPS variance & 0.029 & 0.248 & -- \\
Pareto heavy & fraction=0.05 & DPS AMSE plug-in & 0.050 & -- & -- \\
Pareto heavy & fraction=0.05 & DPS variance & 0.050 & 0.208 & -- \\
Pareto heavy & fraction=0.1 & DPS AMSE plug-in & 0.100 & -- & -- \\
Pareto heavy & fraction=0.1 & DPS variance & 0.100 & 0.156 & -- \\
Pareto light & fraction=0.03 & DPS AMSE plug-in & 0.029 & -- & -- \\
Pareto light & fraction=0.03 & DPS variance & 0.029 & 0.098 & -- \\
Pareto light & fraction=0.05 & DPS AMSE plug-in & 0.050 & -- & -- \\
Pareto light & fraction=0.05 & DPS variance & 0.050 & 0.091 & -- \\
Pareto light & fraction=0.1 & DPS AMSE plug-in & 0.100 & -- & -- \\
Pareto light & fraction=0.1 & DPS variance & 0.100 & 0.083 & -- \\
Abs. Student-3 & fraction=0.03 & DPS AMSE plug-in & 0.029 & 0.116 & 0.895 \\
Abs. Student-3 & fraction=0.03 & DPS variance & 0.029 & 0.116 & 0.894 \\
Abs. Student-3 & fraction=0.05 & DPS AMSE plug-in & 0.050 & 0.146 & 0.917 \\
Abs. Student-3 & fraction=0.05 & DPS variance & 0.050 & 0.146 & 0.917 \\
Abs. Student-3 & fraction=0.1 & DPS AMSE plug-in & 0.100 & 0.285 & 0.905 \\
Abs. Student-3 & fraction=0.1 & DPS variance & 0.100 & 0.285 & 0.905 \\
\bottomrule
\end{tabular}

\endgroup
\end{table}

\begin{table}[htbp]
\centering
\small
\caption{Target sensitivity for \(m=10\), strong allocation and
\(k_j/n_j=0.05\).}
\label{tab:sim:target}
\begingroup
\setlength{\tabcolsep}{4pt}
\begin{tabular}{lrlrrr}
\toprule
DGP & $np$ & Estimator & $\eta_n$ & Log RMSE & Coverage \\
\midrule
Burr heavy & 1 & DPS AMSE plug-in & 0.014 & 0.282 (0.003) & 0.779 (0.006) \\
Burr heavy & 1 & DPS variance & 0.014 & 0.282 (0.003) & 0.779 (0.006) \\
Burr heavy & 5 & DPS AMSE plug-in & 0.051 & 0.221 (0.002) & 0.753 (0.006) \\
Burr heavy & 5 & DPS variance & 0.051 & 0.221 (0.002) & 0.753 (0.006) \\
Burr heavy & 10 & DPS AMSE plug-in & 0.090 & 0.207 (0.002) & 0.697 (0.006) \\
Burr heavy & 10 & DPS variance & 0.090 & 0.207 (0.002) & 0.698 (0.006) \\
Burr light & 1 & DPS AMSE plug-in & 0.359 & 0.076 (0.001) & 0.785 (0.006) \\
Burr light & 1 & DPS variance & 0.359 & 0.076 (0.001) & 0.785 (0.006) \\
Burr light & 5 & DPS AMSE plug-in & 0.724 & 0.069 (0.001) & 0.590 (0.007) \\
Burr light & 5 & DPS variance & 0.724 & 0.069 (0.001) & 0.590 (0.007) \\
Burr light & 10 & DPS AMSE plug-in & 1.014 & 0.075 (0.001) & 0.428 (0.007) \\
Burr light & 10 & DPS variance & 1.014 & 0.075 (0.001) & 0.427 (0.007) \\
Pareto heavy & 1 & DPS AMSE plug-in & 0.014 & -- & -- \\
Pareto heavy & 1 & DPS variance & 0.014 & 0.245 (0.003) & -- \\
Pareto heavy & 5 & DPS AMSE plug-in & 0.051 & -- & -- \\
Pareto heavy & 5 & DPS variance & 0.051 & 0.204 (0.002) & -- \\
Pareto heavy & 10 & DPS AMSE plug-in & 0.090 & -- & -- \\
Pareto heavy & 10 & DPS variance & 0.090 & 0.190 (0.002) & -- \\
Pareto light & 1 & DPS AMSE plug-in & 0.359 & -- & -- \\
Pareto light & 1 & DPS variance & 0.359 & 0.087 (0.001) & -- \\
Pareto light & 5 & DPS AMSE plug-in & 0.724 & -- & -- \\
Pareto light & 5 & DPS variance & 0.724 & 0.090 (0.001) & -- \\
Pareto light & 10 & DPS AMSE plug-in & 1.014 & -- & -- \\
Pareto light & 10 & DPS variance & 1.014 & 0.098 (0.001) & -- \\
Abs. Student-3 & 1 & DPS AMSE plug-in & 0.128 & 0.222 (0.002) & 0.924 (0.004) \\
Abs. Student-3 & 1 & DPS variance & 0.128 & 0.222 (0.002) & 0.924 (0.004) \\
Abs. Student-3 & 5 & DPS AMSE plug-in & 0.297 & 0.142 (0.001) & 0.917 (0.004) \\
Abs. Student-3 & 5 & DPS variance & 0.297 & 0.142 (0.001) & 0.917 (0.004) \\
Abs. Student-3 & 10 & DPS AMSE plug-in & 0.441 & 0.109 (0.001) & 0.880 (0.005) \\
Abs. Student-3 & 10 & DPS variance & 0.441 & 0.109 (0.001) & 0.880 (0.005) \\
\bottomrule
\end{tabular}

\endgroup
\end{table}

\section{Weight stability and limitations}
\label{sec:sim:limitations}

\Cref{tab:sim:weights} reports the AMSE weight diagnostics in the
representative imbalanced design. The oracle and plug-in AMSE weights are
stable in this grid: negative weights are essentially absent, and the average
maximum absolute weight is close to the largest deterministic allocation
weight. The diagnostic matters because it rules out an obvious failure mode:
the lack of AMSE improvement in \Cref{tab:sim:point-main} is not caused by
unstable or frequently negative weights. Rather, in the representative
common-tail design the estimated AMSE rule is numerically close to the simpler
variance allocation, so there is little finite-sample gain to obtain.

The simulation evidence is intentionally limited. It supports the promoted
estimator in the iid common-marginal distributed setting and suggests that
DPS variance weights are a robust practical choice for \(\bm\omega\). It
also shows where the asymptotic theory should not be overinterpreted:
first-order plug-in intervals may under-cover at realistic sample sizes, and
the exact decomposition reveals finite-sample contributions from terms that are
asymptotically negligible under the theorem's scaling. The study does not
address non-identical margins, time series, multivariate targets, random
bridge-Hill weights, or lower-order optimality of \(\bm\nu\). Those questions
would require additional mathematical work before they could become
simulation targets.

\begin{table}[htbp]
\centering
\small
\caption{AMSE weight diagnostics in the representative main design.}
\label{tab:sim:weights}
\begingroup
\setlength{\tabcolsep}{4pt}
\begin{tabular}{lrllrr}
\toprule
DGP & $m$ & Regime & Estimator & Neg. rate & Avg. max $|\omega|$ \\
\midrule
Burr heavy & 10 & strong & DPS AMSE oracle & 0.000 & 0.243 \\
Burr heavy & 10 & strong & DPS AMSE plug-in & 0.000 & 0.243 \\
\bottomrule
\end{tabular}

\endgroup
\end{table}
